\section[Rédaction Bourbaki \redac{397} : divisible et injectif sur un anneau principal]{Rédaction Bourbaki \redac{397} : sur un anneau principal, divisible équivaut à injectif\protect\verifie{Redaction397}}\label{sec:bourbaki-397}

La rédaction \redac{397}, \q{Compléments au chapitre VII d'Algèbre. Dualité des modules de longueur finie sur un anneau principal}, des archives Bourbaki (Archives Henri Poincaré, \ahp{948}) est datée d'avril 1963. Elle définit les modules injectifs, démontre le critère de Baer (sa proposition~2) et en tire un corollaire : sur un anneau principal, un module est injectif si et seulement s'il est divisible ; puis un \q{corollaire du corollaire} : tout quotient d'un module injectif est injectif. Sa transcription reste locale pour l'instant ; on n'en cite ici aucun passage.

\emph{L'idée.} Sur un anneau principal $R$, un idéal est engendré par un seul élément $a$. Prolonger à $R$ une application linéaire $g$ définie sur l'idéal $(a)$, c'est choisir l'image $q$ de $1$, sous la seule contrainte $a\,q = g(a)$, ce qui revient à diviser $g(a)$ par $a$ dans $Q$. Le critère de Baer ramène l'injectivité à ces prolongements depuis les idéaux.

\begin{encadre}
\begin{enonce}[Bourbaki, rédaction \redac{397}, corollaire de la proposition~2 et son corollaire]
Soient $R$ un anneau commutatif et $Q$ un $R$-module.
\begin{enumerate}
\item[a)] Si $R$ est principal et si $Q$ est divisible (pour tout $r \neq 0$, l'application $q \mapsto rq$ de $Q$ dans $Q$ est surjective), alors $Q$ est injectif.
\item[b)] Si $R$ est intègre et si $Q$ est injectif, alors $Q$ est divisible.
\item[c)] Si $R$ est principal et intègre, tout quotient d'un $R$-module injectif est injectif.
\end{enumerate}
\end{enonce}
\end{encadre}

\begin{proof}
a) Par le critère de Baer, il suffit de prolonger à $R$ toute application linéaire $g$ définie sur un idéal $(a)$. Si $a = 0$, on prend l'application nulle. Sinon, on choisit $q$ tel que $a q = g(a)$ ; l'application $r \mapsto r q$ prolonge $g$, car $g(c a) = c\,g(a) = (c a)\, q$.

b) Soient $r \neq 0$ et $q \in Q$. Comme $R$ est intègre, $c \mapsto c r$ est une bijection de $R$ sur l'idéal $(r)$, et $c r \mapsto c q$ est une application linéaire de $(r)$ dans $Q$. Elle se prolonge à $R$ ; si $x$ est l'image de $1$, alors $r x$ est l'image de $r$, c'est-à-dire $q$.

c) Si $Q$ est injectif, il est divisible par b), et un quotient d'un module divisible est divisible ; il est donc injectif par a).
\end{proof}

\noindent Le point a) n'utilise pas l'intégrité : il vaut sur tout anneau commutatif principal.

\begin{sloppypar}
\noindent\emph{Vérification en Lean.} a) est \texttt{baer\_of\_divisible} et \texttt{injective\_of\_divisible}, b) \texttt{divisible\_of\_injective}, l'équivalence \texttt{injective\_iff\_divisible}, et c) \texttt{injective\_quotient}, dans \href{\leanblobmain/Redaction397.lean}{\texttt{Redaction397.\allowbreak{}lean}}. Le critère de Baer de mathlib demande que $R$ et $Q$ soient dans le même univers, et la formalisation garde cette restriction.
\end{sloppypar}

\litterature{retrouvée. Le critère est de Baer (1940) ; l'équivalence entre divisible et injectif sur un anneau principal est classique (Lam, \emph{Lectures on Modules and Rings}, §~3). mathlib n'en avait que le cas de $\mathbf Z$ (\texttt{Module.Baer.of\_divisible}) ; le cas général est proposé pour mathlib, et l'édition ne le donne pas pour nouveau.}
